\section{Conclusão}

\begin{frame}{Limitações e trabalhos futuros}
	\begin{itemize}
		\item Dois otimizadores Adam (rede siamesa e modelo de velocidade) precisam de taxas de aprendizado bem coordenadas entre si.
		\item Sem camada de \textit{pooling}, o tamanho espacial é preservado em toda a rede: em dados de campo de grande dimensão, isso aumenta o custo de memória e processamento.
		\item A rede siamesa aprendida em um conjunto de dados não transfere bem para outro: \textit{transfer learning} direto do Overthrust para o Marmousi2 piora o resultado (SNR cai de 22,57 para 20,27 dB).
	\end{itemize}
\end{frame}

\begin{frame}{Conclusão}
	\begin{itemize}
		\item O SiameseFWI usa uma rede siamesa \textit{self-supervised}, treinada durante o próprio FWI, para comparar dado observado e simulado em um espaço latente.
		\item Supera FWI com perdas convencionais (correlação cruzada, Wasserstein) e o FWIGAN, tanto em modelos sintéticos quanto em dado de campo real.
		\item Alcança esse ganho com apenas 10\%--20\% de custo computacional extra, e se mantém robusto a ruído e a modelo inicial impreciso.
	\end{itemize}
\end{frame}
